\chapter{La bicouche : information paire et transport géométrique impair}
\label{chap:paridad-bicapa-barbero-rev6}
\index{entropie!bicouche}
\index{Barbero--Immirzi!parité}
\index{feuillet!échange}

Barbero--Immirzi et l'entropie des deux feuillets ne sont pas deux quantités
juxtaposées. Ils naissent des mêmes canaux \(q_\pm\), mais conservent des données
distinctes. La normalisation probabiliste oublie l'échelle commune et mesure la
répartition entre orientations ; la fonctionnelle hexade--octade conserve le
signe de cette orientation. La première lecture est paire et la seconde impaire.

\section{Distribution normalisée entre feuillets}

Écrivons
\[
 x=A_{\rm rad},
 \qquad
 y=C^*_{\rm rad},
 \qquad
 q_-=e^{-x+y},
 \qquad
 q_+=e^{-x-y}.
\]
La somme \(q_-+q_+=2e^{-x}\cosh y\) permet de définir
\[
 p_-=\frac{q_-}{q_-+q_+}
     =\frac{e^y}{2\cosh y},
 \qquad
 p_+=\frac{q_+}{q_-+q_+}
     =\frac{e^{-y}}{2\cosh y}.
\]
Le facteur commun \(e^{-x}\) disparaît. La distribution conserve l'asymétrie
entre feuillets, mais non l'échelle angulaire commune.

\begin{definition}[Entropie bicouche]
\label{def:entropia-bicapa-rev6}
L'entropie de la distribution orientée est
\[
 S_{\rm bi}(y)
 :=
 -p_-\log p_--p_+\log p_+
 =
 \log(2\cosh y)-y\tanh y.
\]
\end{definition}

\begin{proposition}[Parité et maximum de l'entropie bicouche]
\label{prop:paridad-entropia-bicapa-rev6}
La fonction \(S_{\rm bi}\) est paire et satisfait
\[
 S_{\rm bi}(-y)=S_{\rm bi}(y),
 \qquad
 S_{\rm bi}'(y)=-y\,\operatorname{sech}^2y,
 \qquad
 0<S_{\rm bi}(y)\leq\log2.
\]
Le maximum \(\log2\) est atteint uniquement en \(y=0\), lorsque les deux feuillets
ont le même poids.
\end{proposition}

\begin{proof}
La parité découle du fait que \(\cosh\) est paire et que \(y\tanh y\) l'est également.
La dérivation directe produit
\[
 \frac{\dd}{\dd y}\log(2\cosh y)=\tanh y,
 \qquad
 \frac{\dd}{\dd y}(y\tanh y)
 =\tanh y+y\operatorname{sech}^2y.
\]
Par conséquent \(S_{\rm bi}'(y)=-y\operatorname{sech}^2y\) : la fonction croît sur
\((-\infty,0)\), décroît sur \((0,\infty)\) et prend à l'origine la valeur
\(\log2\). La positivité est celle de l'entropie d'une distribution à deux
poids strictement positifs.
\end{proof}

\section{Parité du transport hexade--octade}

Soit
\[
 s_n(q_-,q_+)=q_-^n-q_+^n.
\]
L'échange des feuillets est l'involution
\[
 \mathsf J_{\rm sh}:(x,y)\longmapsto(x,-y),
\]
qui permute \(q_-\leftrightarrow q_+\). D'où,
\[
 s_n\circ\mathsf J_{\rm sh}=-s_n.
\]
La fonctionnelle construite dans le \cref{mdg:chap:barbero} est
\[
 \Gamma_{6\to8}(A,C^*)
 =
 \frac{\pi A C^*}{180}
{}+12\bigl(S_{90}(q_-)-S_{90}(q_+)\bigr)
-\bigl(S_{120}(q_-)-S_{120}(q_+)\bigr).
\]

\begin{theorem}[Théorème de parité bicouche]
\label{thm:paridad-bicapa-barbero-rev6}
Sous l'échange \(C^*\mapsto-C^*\),
\[
 \boxed{S_{\rm bi}(-C^*_{\rm rad})
        =S_{\rm bi}(C^*_{\rm rad}),}
\]
\[
 \boxed{\Gamma_{6\to8}(A,-C^*)
        =-\Gamma_{6\to8}(A,C^*).}
\]
La même bicouche produit donc une lecture informationnelle paire et un
transport géométrique impair.
\end{theorem}

\begin{proof}
La première égalité est la
\cref{prop:paridad-entropia-bicapa-rev6}. Dans la seconde, le terme
\(\pi AC^*/180\) change de signe et chaque différence de Lambert change de
signe parce que ses deux arguments sont échangés. La combinaison complète est
impaire.
\end{proof}

\section{Interprétation de l'opérateur constitutif comme couplage entre vide, champ et réponse}

L'entropie bicouche et \(\Gamma_{6\to8}\) ne sont pas deux noms pour le même
nombre. La première quantifie le poids que partagent les deux orientations
après annulation de l'échelle commune. Le second conserve l'orientation et
l'incidence du passage \(6\to8\). Au point HMT,
\[
 S_{\rm bi}
 =0.69246069960358548\ldots\ {\rm nats},
 \qquad
 \Gamma_{6\to8}
 =0.27400767269445927\ldots.
\]
Le lien entre information et géométrie consiste en ce que les deux lectures
se factorisent par les mêmes canaux \(q_\pm\), et non à égaler leurs valeurs.

La bande particule--antiparticule du
\cref{chap:banda-particula-virtual} reproduit la même distinction de
parité : la masse et la quantité de corrélation descendent au quotient de
bande, tandis que charge, orientation et transport changent de signe. Dans
l'action de Cartan--Holst, \(\Gamma_{6\to8}\) occupe la coordonnée orientée ;
dans la lecture informationnelle, \(S_{\rm bi}\) compte la distribution entre
feuillets. Cela explique comment une seule structure peut conserver
l'information et inverser l'orientation sans confondre les deux opérations.
